\chapter{State-anchored one-form bimodule: каноническая CP-map}

Inner-fluctuation gate оставил выбор: малая family algebra не содержит
CP-capable correction, а полная \(M_3(\mathbb C)\) возвращает девять
произвольных directions. Вместо расширения algebra выделим минимальный
one-form module, определённый уже найденным pure state \(\rho_\star\).

\section{State-supported module}

Положим \(Q_\star=I-\rho_\star\) и потребуем, чтобы каждый subleading
connector касался selected state хотя бы одним endpoint. Тогда допустимое
пространство равно
\begin{equation}
 \mathcal M_{\rho}:=
 \rho_\star M_3+M_3\rho_\star
 =\{X\in M_3:Q_\star XQ_\star=0\}.
\end{equation}
Это не новая полная algebra, а *-замкнутый bimodule над state algebra
\begin{equation}
 \mathcal A_\rho=
 \operatorname{span}\{\rho_\star,Q_\star\}.
\end{equation}
Прямой rank test даёт
\begin{equation}
 \dim_{\mathbb C}\mathcal M_\rho=5,
 \qquad
 \dim_{\mathbb C}\ker\Pi_\rho=4.
\end{equation}
Исключённые четыре directions образуют внутренний block
\(Q_\star M_3Q_\star\), не связанный с выбранным state.

\section{Единственная trace-orthogonal projection}

Относительно уже используемой Hilbert--Schmidt trace metric ортогональная
projection на \(\mathcal M_\rho\) единственна:
\begin{equation}
 \boxed{
 \Pi_\rho(X)=X-Q_\star XQ_\star
 =\rho_\star X+X\rho_\star-\rho_\star X\rho_\star.}
\end{equation}
Она idempotent, self-adjoint как superoperator, *-compatible и covariant:
\begin{equation}
 \Pi_{U\rho U^\dagger}(UXU^\dagger)
 =U\Pi_\rho(X)U^\dagger.
\end{equation}
Машинные residuals всех этих identities не превосходят
\(5\times10^{-13}\).

Тем самым support principle и существующая trace metric выбирают map без
continuous coefficient:
\begin{equation}
 Y_s^{\rm anch}=P_-+i\Pi_{\rho_\star}(H_s),
 \qquad s=u,d.
\end{equation}

\section{Первый выбранный CP readout}

Полученные normalized spectra равны
\begin{align}
 m_u/m_{u,\max}&=(0.192658,\,0.375569,\,1),\\
 m_d/m_{d,\max}&=(0.119012,\,0.430519,\,1),
\end{align}
а cubic CP invariant имеет значение
\begin{equation}
 \Im\Tr[M_u,M_d]^3=-0.287960\ne0.
\end{equation}
Absolute mixing matrix имеет вид
\begin{equation}
 |V_{\rm anch}|\simeq
 \begin{pmatrix}
 0.4302&0.8512&0.3007\\
 0.8349&0.4751&0.2777\\
 0.3432&0.2230&0.9124
 \end{pmatrix}.
\end{equation}
Следовательно, map является genuinely CP-capable, но этот минимальный ledger
ещё не проходит observed small-mixing и mass-hierarchy scorecard.

\begin{theorem}[Conditional state-anchored selector]
После выбора \(\rho_\star\) вариационным family functional support axiom
\(Q_\star XQ_\star=0\) и Hilbert--Schmidt metric однозначно выделяют
пятимерный bimodule и projection \(\Pi_\rho\). Полученная Yukawa map не имеет
continuous flavour coefficient и порождает ненулевой CP invariant.
\end{theorem}

\begin{remark}
Это сильнее предыдущего existence pass: Jordan и произвольные full-\(M_3\)
maps больше не равноправны, поскольку только \(\Pi_\rho\) является
ортогональной projection на state-supported module. Условным остаётся сам
support axiom; его ещё нужно вывести из graph incidence или parent action.
\end{remark}

\begin{nogocriterion}[Запрет объявления phenomenological pass]
Ненулевой CP invariant не компенсирует крупное mixing и слабую hierarchy.
До появления четырёх sector-sensitive edges и RG readout найденная map
считается структурным, а не феноменологическим pass.
\end{nogocriterion}

Следующая задача сужена: вывести support axiom из одного finite graph и
различить \(u,d,e,\nu\) без четырёх независимых maps. Машинный ledger сохранён
в \path{s2t_v4_state_anchored_bimodule_gate_results.json}.